\documentclass{article}

% Pass the round citation option to natbib BEFORE the NeurIPS style auto-loads it.
% The style loads natbib itself, so a second \usepackage[round]{natbib} would cause
% an "Option clash for package natbib" error; passing the option here avoids it.
\PassOptionsToPackage{round}{natbib}

% Keep the same paper style when the NeurIPS style file is available.
\IfFileExists{neurips_2026.sty}{\usepackage[preprint]{neurips_2026}}{}

\usepackage[utf8]{inputenc}
\usepackage[T1]{fontenc}
\usepackage{url}
\usepackage{booktabs}
\usepackage{amsfonts}
\usepackage{microtype}
\usepackage{amsmath}
\usepackage{amssymb}
\usepackage{graphicx}
\usepackage{float}
\usepackage{enumitem}
% natbib is loaded by the NeurIPS style above (with the round option passed via
% \PassOptionsToPackage near \documentclass); do not load it again here.
\usepackage[hidelinks,colorlinks=true,citecolor=blue,linkcolor=blue,urlcolor=blue]{hyperref}
\usepackage[capitalize,noabbrev]{cleveref}

% >>> BEGIN COPY INTO A DOCUMENT PREAMBLE >>>
% Copied verbatim from the shared skill's style reference
% /p/rlprojects/.claude/skills/rnd-experiment-tex-track/example.tex (its marked block).
% WARNING: \arrayrulecolor is a GLOBAL switch — only use it through the two rule macros below,
% which both end with the color on black. A \multicolumn that is NOT the last cell of its row
% must end its spec with ? (e.g. {c?}), or the grey vertical on its right edge disappears.
\usepackage[table]{xcolor}  % [table] loads colortbl -> \arrayrulecolor
\usepackage{array}
\setlength{\arrayrulewidth}{0.5pt}                 % \hline/\vline only; booktabs rules unaffected
\definecolor{tabgrey}{gray}{0.65}
\newcommand{\greyvline}{{\color{tabgrey}\vline}}   % braces keep the color inside the rule
\newcolumntype{?}{!{\greyvline}}                   % column separator: one grey vertical rule
\newcommand{\greyrule}{\arrayrulecolor{tabgrey}\hline\arrayrulecolor{black}} % between rows of one category
\newcommand{\blackrule}{\arrayrulecolor{black}\hline}                        % category boundary / frame
% <<< END COPY <<<

% Visible-edit markup of the shared skill claude-edit-latex: a replacement of author text is
% \sout{old}\claudeedit{new} (struck original + brown replacement), and a rationale annotation is
% \claudecomment{...} (dark green). Both macros stay defined across every acceptance stage.
\usepackage[normalem]{ulem}
\definecolor{claude-brown}{RGB}{133,45,15}        % brown --- replacement-text color
\definecolor{claude-green}{RGB}{0,100,0}          % dark green --- annotation color
\newcommand{\claudeedit}[1]{{\color{claude-brown}\,#1}}
\newcommand{\claudecomment}[1]{({\color{claude-green} {Claude: #1}})}

% The NeurIPS style sets \flushbottom, which stretches a page's vertical spacing until the text
% reaches the bottom margin. This document's pages are mostly one tall table each, so a page that
% a table did not fit on gets its few paragraphs spread across the whole height with large gaps
% between them. \raggedbottom leaves that space at the bottom instead.
\raggedbottom

\title{Exploration Bonuses on the JAX Training Platform:\\ Environment and Train-Run Development Document}

\author{%
  Anonymous Author(s)\\
  Affiliation\\
  \texttt{email}
}

% Tables and figures are numbered by section (Table 1.1, Figure 1.2, ...), so a reader who is
% told "the table in 1.1" can find it, and so a later section's artifacts never renumber this
% section's. amsmath, loaded above, provides \numberwithin.
\numberwithin{table}{section}
\numberwithin{figure}{section}

\begin{document}

\maketitle

\tableofcontents
\clearpage

%==============================================================================
\section{PointMaze on the JAX platform}
\label{sec:pm-jax}
% The environment is subsection 1.0 and the first train run is 1.1, so a later run on the same
% environment becomes 1.2 and the shared environment specification keeps the .0 slot.
\setcounter{subsection}{-1}

\subsection{Environment specification}
\label{sec:pm-jax-envspec}

The platform trains on one environment: the large PointMaze with the reset noise switched off.
Every episode begins at exactly the same point and aims at exactly the same goal, so two copies
that differ only by their policy sampling see the same task, and a difference between two
configurations is not a difference between two randomly drawn start positions.
\Cref{tab:pm-jax-env} states it knob by knob.

\begin{table}[H]
\centering
\footnotesize
\setlength{\tabcolsep}{3pt}
\renewcommand{\arraystretch}{1.2}
\resizebox{\textwidth}{!}{%
\begin{tabular}{@{}>{\raggedright\arraybackslash}p{3.4cm}
?>{\raggedright\arraybackslash}p{6.4cm}
?>{\raggedright\arraybackslash}p{4.4cm}@{}}
\blackrule
\textbf{Parameter} & \textbf{PointMaze large, no reset noise} & \textbf{Where the value is set} \\
& \small(the platform's only environment so far) & \small(field of the copied code) \\
\blackrule
\multicolumn{3}{@{}l}{\textit{Identity and provenance}} \\
\greyrule
specification name & \texttt{pointmaze\allowbreak\_large\allowbreak\_cont400\allowbreak\_nonoise@1} &
  the name a run manifest carries \\
\greyrule
provenance & extracted from \texttt{09\allowbreak\_parallelization} at tag
  \texttt{jax-rnd-baseline-v0.1.0}; nothing but import paths changed on the copy &
  \texttt{ORIGIN.md} \\
\greyrule
implementation & \texttt{envs/pointmaze/\allowbreak jax\allowbreak\_pointmaze.py} (the stepper) and
  \texttt{pm\allowbreak\_common.py} (physics constants, maps, configuration) &
  \texttt{src/exploration\allowbreak\_platform/envs/\allowbreak pointmaze/} \\
\blackrule
\multicolumn{3}{@{}l}{\textit{Maze and spaces}} \\
\greyrule
maze map & \texttt{large}: 12 columns $\times$ 9 rows of 1\,m cells, walls on the border,
  46 open cells & \texttt{EnvConfig.\allowbreak map\allowbreak\_name} \\
\greyrule
world extent & $x \in [-6, 6]$, $y \in [-4.5, 4.5]$; the origin is the maze centre and row $0$ is
  the top row & \texttt{pm\allowbreak\_common.\allowbreak cell\allowbreak\_center} \\
\greyrule
agent body & 2-D ball, force-actuated, radius $0.1$\,m & \texttt{pm\allowbreak\_common} physics
  constants \\
\greyrule
observation & $4$-d: $x$, $y$, $v_x$, $v_y$. The goal is \emph{not} observed &
  \texttt{jax\allowbreak\_pointmaze.\allowbreak step} \\
\greyrule
action & $2$-d force, each component clipped to $[-1, 1]$ &
  \texttt{pm\allowbreak\_common.\allowbreak ACT\allowbreak\_CLIP} \\
\blackrule
\multicolumn{3}{@{}l}{\textit{Start and goal}} \\
\greyrule
start cell & row $7$, column $1$ (row $0$ is the top); world position $(-4.5, -3.0)$ &
  \texttt{EnvConfig.\allowbreak start\allowbreak\_cell} \\
\greyrule
goal cell & row $1$, column $10$; world position $(4.5, 3.0)$ &
  \texttt{EnvConfig.\allowbreak goal\allowbreak\_cell} \\
\greyrule
reset position noise & $0$. Every episode starts exactly at the start cell's centre and the goal
  sits exactly at its own centre; the only randomness left in an episode is the policy's own
  sampling & \texttt{EnvConfig.\allowbreak position\allowbreak\_noise} \\
\greyrule
initial velocity & $(0, 0)$ at every reset &
  \texttt{jax\allowbreak\_pointmaze.\allowbreak \_spawn} \\
\greyrule
goal radius & $0.45$\,m: a step whose distance to the goal is at most this earns reward $1$ &
  \texttt{EnvConfig.\allowbreak goal\allowbreak\_radius} \\
\blackrule
\multicolumn{3}{@{}l}{\textit{Episode and reward}} \\
\greyrule
episode length cap & $400$ steps; one step is one $0.01$\,s integrator step &
  \texttt{EnvConfig.\allowbreak max\allowbreak\_episode\allowbreak\_steps} \\
\greyrule
task type & continuing: an episode is never terminated at the goal, only truncated at the cap, so
  a copy that reaches the goal keeps collecting reward until the cap &
  \texttt{EnvConfig.\allowbreak continuing\allowbreak\_task} \\
\greyrule
reward & sparse: $1$ on a step inside the goal radius, $0$ otherwise; no shift is added &
  \texttt{EnvConfig.\allowbreak reward\allowbreak\_shift} \\
\greyrule
discount, extrinsic & $0.999$ & \texttt{PPOConfig.\allowbreak gamma\allowbreak\_ext} \\
\greyrule
discount, intrinsic & $0.99$ (used only by an arm that has an intrinsic bonus) &
  \texttt{PPOConfig.\allowbreak gamma\allowbreak\_int} \\
\blackrule
\end{tabular}%
}
\caption{The environment every train run in this section uses. The two discounts are the trainer's
defaults rather than environment fields, and are listed here because they define the objective the
environment is optimised against. Long typewriter names break after every escaped underscore
(\texttt{\textbackslash allowbreak}).}
\label{tab:pm-jax-env}
\end{table}

\subsection{Train run 1.1: the first algorithm batch}
\label{sec:pm-jax-run11}

Four algorithm arms run at once: random network distillation, no exploration bonus at all, and the
two oracle visit-count bonuses that differ only in how fast their bonus falls with the count. Each
arm is swept over three learning rates, and each arm that has a bonus is swept over eleven weights
on that bonus.

\paragraph{Hypotheses.}
\begin{enumerate}[label=\textbf{H\arabic*},leftmargin=*,labelsep=0.6em]
\item \textbf{Random network distillation beats no exploration bonus.} At every one of the three
  learning rates, \texttt{rnd\allowbreak\_next\allowbreak\_state} at its best intrinsic weight
  reaches a higher maze-cell coverage and a higher whole-run reward than
  \texttt{no\allowbreak\_exploration} at the same learning rate.
\item \textbf{The oracle bonuses bound random network distillation from above.} The visit-count
  arms are given the true novelty signal that random network distillation has to estimate from a
  predictor's error, so at their best intrinsic weight they reach at least the coverage and reward
  that the best random-network-distillation configuration reaches.
\item \textbf{The two decays trade coverage against reward differently.} $1/\sqrt{n}$ keeps paying
  for a state long after $1/n$ has stopped, so the $1/\sqrt{n}$ arm reaches the higher coverage and
  the $1/n$ arm reaches the higher reward once the goal has been found.
\item \textbf{Reward against intrinsic weight is single-peaked.} For every bonus arm, whole-run
  reward as a function of the intrinsic weight rises and then falls: a weight far below the peak
  leaves the agent exploring too little to find the goal, and a weight far above it drowns the
  sparse extrinsic reward in the bonus, so the agent keeps chasing novelty after the goal is known.
\end{enumerate}

\begin{table}[H]
\centering
\footnotesize
\setlength{\tabcolsep}{3pt}
\renewcommand{\arraystretch}{1.2}
\resizebox{\textwidth}{!}{%
% Widths: the arm columns are sized by their content, and every spanning cell's width is the sum
% of the columns it covers (11.3 = 3.1+2.3+2.95+2.95, 5.9 = 2.95+2.95).
\begin{tabular}{@{}>{\raggedright\arraybackslash}p{2.9cm}
?>{\raggedright\arraybackslash}p{3.1cm}
?>{\raggedright\arraybackslash}p{2.3cm}
?>{\raggedright\arraybackslash}p{2.95cm}
?>{\raggedright\arraybackslash}p{2.95cm}@{}}
\blackrule
\textbf{Hyperparameter} & \textbf{\texttt{rnd\allowbreak\_next\allowbreak\_state}} &
\textbf{\texttt{no\allowbreak\_exploration}} &
\textbf{\texttt{gt\allowbreak\_position\allowbreak\_velocity\allowbreak\_sqrt}} &
\textbf{\texttt{gt\allowbreak\_position\allowbreak\_velocity\allowbreak\_linear}} \\
& \small(random network distillation) & \small(no intrinsic bonus) &
\small(oracle counts, $1/\sqrt{n}$) & \small(oracle counts, $1/n$) \\
\blackrule
\multicolumn{5}{@{}l}{\textit{Environment and RL algorithm}} \\
\greyrule
environment specification & \multicolumn{4}{>{\raggedright\arraybackslash}p{11.3cm}}{the
  no-noise large PointMaze of \Cref{tab:pm-jax-env}; this run adds no new
  environment-specification table} \\
\greyrule
base RL algorithm & proximal policy optimisation, this platform's JAX trainer &
  \multicolumn{3}{c}{same} \\
\greyrule
update style & one gradient step per rollout (\texttt{full\allowbreak\_batch}) &
  \multicolumn{3}{c}{same} \\
\greyrule
rollout length $T$ & $128$ steps & \multicolumn{3}{c}{same} \\
\greyrule
environments per copy $N$ & $4$, so one iteration collects $T \times N = 512$ transitions per copy &
  \multicolumn{3}{c}{same} \\
\greyrule
environment steps per copy & \sout{$10{,}035{,}200$: $19{,}600$ iterations of $512$ steps. That is
  $10^{7}$ to within $0.35$ per cent, and unlike $10^{7}$ it is a whole number of $200$-iteration
  recording windows}\claudeedit{$10$M} & \multicolumn{3}{c}{same} \\
\blackrule
\multicolumn{5}{@{}l}{\textit{Intrinsic bonus}} \\
\greyrule
bonus signal & the prediction error of a trained network against a fixed random target & N/A &
  \multicolumn{2}{>{\centering\arraybackslash}p{5.9cm}}{the oracle count of the discrete next
  state} \\
\greyrule
bonus input & the next state $(x, y, v_x, v_y)$ & N/A &
  \multicolumn{2}{>{\centering\arraybackslash}p{5.9cm}}{the next state's maze cell and its two
  velocity bins} \\
\greyrule
bonus value & the squared prediction error, whitened by its running standard deviation & N/A &
  $\min(1, n^{-1/2})$ at count~$n$ & $\min(1, n^{-1})$ at count~$n$ \\
\greyrule
what the bonus holds & a predictor network trained with the agent: 2 hidden layers of $256$, $128$
  outputs & nothing & \multicolumn{2}{>{\centering\arraybackslash}p{5.9cm}}{one integer count table
  per copy, $9 \times 12 \times 100 = 10{,}800$ entries; a wall or unvisited cell scores $1$} \\
\greyrule
when the bonus updates & every gradient step, as a term of the joint loss & N/A &
  \multicolumn{2}{>{\centering\arraybackslash}p{5.9cm}}{once per rollout: one scatter that counts
  it, one gather that scores it} \\
\greyrule
where the bonus enters & the advantage: $\beta \times$ intrinsic $+\ 2 \times$ extrinsic & N/A &
  \multicolumn{2}{>{\centering\arraybackslash}p{5.9cm}}{the advantage: $\beta \times$ intrinsic
  $+\ 2 \times$ extrinsic} \\
\blackrule
\multicolumn{5}{@{}l}{\textit{Swept knobs}} \\
\greyrule
learning rate & \multicolumn{4}{c}{$10^{-3}$, $10^{-4}$, $10^{-5}$ --- 3 values} \\
\greyrule
intrinsic weight $\beta$ & $10^{-5}$ to $10^{5}$ in factor-of-ten steps --- 11 values & N/A &
  \multicolumn{2}{>{\centering\arraybackslash}p{5.9cm}}{$10^{-5}$ to $10^{5}$ in factor-of-ten
  steps --- 11 values} \\
\blackrule
\multicolumn{5}{@{}l}{\textit{Run bookkeeping (sweep config counts)}} \\
\greyrule
configs in this column & $3$(learning rate) $\times$ $11$(intrinsic weight) $=33$ &
  $3$(learning rate) $=3$ & $3$(learning rate) $\times$ $11$(intrinsic weight) $=33$ &
  $3$(learning rate) $\times$ $11$(intrinsic weight) $=33$ \\
\greyrule
total configs in this table & \multicolumn{4}{c}{$33+3+33+33=\mathbf{102}$} \\
\greyrule
copies in this column & $33 \times 256 = 8{,}448$ & $3 \times 256 = 768$ &
  $33 \times 256 = 8{,}448$ & $33 \times 256 = 8{,}448$ \\
\greyrule
seeds & \multicolumn{4}{>{\raggedright\arraybackslash}p{11.3cm}}{$256$ copies per configuration,
  seeded in pairs: copy $k$ of every configuration starts from the same weights and meets the same
  environment draws, so a difference between two configurations is the swept values' doing and not
  a difference between seeds} \\
\blackrule
\multicolumn{5}{@{}l}{\textit{Extension}} \\
\greyrule
extended at a larger scale & \multicolumn{4}{>{\raggedright\arraybackslash}p{11.3cm}}{\sout{the
  best configuration of each of these four arms --- every one of them at the learning rate
  $10^{-3}$, at $\beta = 10^{1}$ for \texttt{rnd\allowbreak\_next\allowbreak\_state} and
  $\beta = 10^{0}$ for the two oracle arms --- was re-run at $1{,}024$ copies for
  $1{,}000{,}038{,}400$ environment steps per copy, that is $1{,}953{,}200$ iterations and $9{,}766$
  recording windows: $4$ times the copies of one cell above and about $100$ times its step
  budget.}\claudeedit{each arm's best configuration --- ranked at $10$M steps off curves that were
  still rising --- re-run at $1{,}024$ copies for $1000$M steps per copy; results below the first
  double rule of \Cref{tab:pm-jax-run11-results}}} \\
\greyrule
\claudeedit{extended a second time} &
  \multicolumn{4}{>{\raggedright\arraybackslash}p{11.3cm}}{\claudeedit{the whole grid of the two
  arms that were still improving at $1000$M steps ---
  \texttt{rnd\allowbreak\_next\allowbreak\_state} and
  \texttt{gt\allowbreak\_position\allowbreak\_velocity\allowbreak\_sqrt}, $33$ configurations each
  --- re-run at $128$ copies per configuration for $1000$M steps per copy, to see whether the
  configuration a $10$M-step sweep picks is still the right one at $1000$M; results below the
  second double rule of \Cref{tab:pm-jax-run11-results}}} \\
\blackrule
\end{tabular}%
}
\caption{What train run 1.1 sweeps, and what its \sout{extension}\claudeedit{two extensions}
re-ran. ``same'' collapses a value
shared with the columns to its
left; the grey verticals make each span's extent visible. ``N/A'' marks a knob that does not apply
--- \texttt{no\allowbreak\_exploration} has no bonus, so it has no bonus knobs and no intrinsic
weight to sweep. The two oracle columns repeat their shared values rather than writing ``same'',
because the column to their left is the one arm the value does not apply to. The last
\sout{row is the extension's}\claudeedit{two rows are the two extensions'}; every row above
\sout{it}\claudeedit{them} describes the original sweep and is unchanged from its launch.}
\label{tab:pm-jax-run11-swept}
\end{table}

\noindent\claudecomment{Two edits in \Cref{tab:pm-jax-run11-swept}, both under the shared skill
\texttt{rl-experiment-writing}: the step-budget row now reads $10$M, and the extension row is cut
from $118$ words to $24$. Every exact figure the two rows carried --- $10{,}035{,}200$ steps per
copy over $19{,}600$ iterations for the sweep, $1{,}000{,}038{,}400$ over $1{,}953{,}200$ for the
extension --- is recorded in the runs' own \texttt{manifest.yaml} files
(\texttt{compile\allowbreak\_signature}, and \texttt{extends.\allowbreak
configurations\allowbreak\_extended} in the extension's), which is where a reader checking
reproducibility looks. The extension's $200$-character run-folder path was deleted outright rather
than struck through, because rule~2 of that skill keeps run-folder paths out of the document
entirely; the run is named by date in the status paragraph below and in the caption of
\Cref{tab:pm-jax-run11-results}.}

\noindent\claudecomment{One row added to \Cref{tab:pm-jax-run11-swept} for the second extension of
2026-08-17, written under the shared skill \texttt{rl-experiment-writing}: which two arms, from
what to what, and why, in $46$ words and with no run-folder path. Its exact budget ---
$1{,}000{,}038{,}400$ environment steps per copy over $1{,}953{,}200$ iterations, $128$ copies of
each of $33$ configurations per arm --- is in that run's \texttt{manifest.yaml}
(\texttt{compile\allowbreak\_signature} and \texttt{extends.\allowbreak arms\allowbreak\_extended}).
The row above it now says ``the first double rule'' rather than ``the double rule'', because
\Cref{tab:pm-jax-run11-results} has two.}

\paragraph{Results --- status.}
The run is complete. All four work units finished on 2026-08-16 between 21:04 and 21:36 PT: the two
oracle visit-count units and the distillation unit on three H100 NVL cards of the scheduler's
\texttt{gpu} partition, and the unit with no bonus on the unscheduled H100 machine
\texttt{serval05}. Every unit wrote its $98$ recorded windows and its completion record, so all
$102$ configurations of the sweep --- $26{,}112$ copies --- are scored, and every window used is
phase-blocked. \Cref{tab:pm-jax-run11-results} and \Cref{fig:pm-jax-run11-curves} are written by the
run's generator
(\texttt{development\allowbreak\_document/code/\allowbreak 2026-08-16-21-00\allowbreak\_pm-jax-run11-tables/\allowbreak make\allowbreak\_all.py})
and by \texttt{analysis/\allowbreak code/\allowbreak make\allowbreak\_reward\allowbreak\_curves.py}
in the run folder; both import the run's own aggregation module,
\texttt{code/\allowbreak aggregate.py}, rather than re-deriving any number, so the two cannot name
different winners.

\textbf{The extension is also complete.} The four winning configurations above were re-run at
$1{,}024$ copies for \sout{$1{,}000{,}038{,}400$ environment}\claudeedit{$1000$M} steps per copy ---
a hundred times the budget
--- on 2026-08-16 between 23:19 PT and 2026-08-17 03:44 PT, on five H100 NVL cards of the
\texttt{gpu} partition. The distillation arm ran as two chunks of $512$ copies holding copy indices
$0$--$511$ and $512$--$1023$; the other three ran as one chunk of $1{,}024$ each. Every chunk wrote
\sout{its $9{,}766$ recorded windows}\claudeedit{all of its recorded windows} and its completion
record, all $48{,}830$ windows are phase-blocked,
and the chunks' copy-index ranges were checked to partition each arm's $1{,}024$ copies exactly
once, so every extension row is a mean over the full $1{,}024$. \claudeedit{Its} rows are the ones
below
the double rule of \Cref{tab:pm-jax-run11-results}. The figure now shows the extension's curves;
which cards each chunk went to, and what every rejected alternative would have cost, is
\texttt{code/\allowbreak submission\allowbreak\_plan.json} in \sout{that}\claudeedit{the extension
run} folder.
\claudecomment{Rounding and path removal under the shared skill \texttt{rl-experiment-writing}:
the exact budget ($1{,}000{,}038{,}400$ steps per copy, $1{,}953{,}200$ iterations, $9{,}766$
windows per chunk) is recorded in the extension run's \texttt{manifest.yaml}, section
\texttt{compile\allowbreak\_signature}, and repeated in its
\texttt{experiment\allowbreak\_background.md}. The sentence naming the run folder was deleted
outright rather than struck through, since rule~2 of that skill keeps such paths out of the
document entirely; the run is identified here by its dates and in the caption of
\Cref{tab:pm-jax-run11-results}. The counts kept exact are the ones whose exactness is the claim
--- that all $48{,}830$ windows are phase-blocked and that the chunks partition the $1{,}024$
copies once each.}

\claudeedit{\textbf{The second extension is complete.} The whole grid of the two arms that were
still improving --- random network distillation and the oracle $1/\sqrt{n}$ bonus, $33$
configurations each at $128$ copies --- was re-run at $1000$M steps per copy on 2026-08-17 and
2026-08-18, because the first extension had tested only the configuration each arm won with at
$10$M steps. All $66$ configurations are scored over the full $128$ copies; its rows are the ones
below the second double rule of \Cref{tab:pm-jax-run11-results}.}
\claudecomment{Written under the shared skill \texttt{rl-experiment-writing}: which
configurations, from what to what, and why, with no run-folder path and no exact step count. The
run is named by its dates, as the first extension is. Its work was cut into $44$ chunks over $16$
node classes and one of its cards failed mid-run, which is recorded in the run's
\texttt{infra\allowbreak\_history.md} and deliberately not here --- rule~3 of that skill keeps
infrastructure telemetry out of this document.}

\paragraph{How to read the metrics tables.}
Define a \emph{configuration} as one (learning rate, intrinsic weight) cell of one algorithm arm,
and define $C$ as the number of copies it is given, \sout{$C = 256$ here}\claudeedit{$C = 256$
in the original sweep, $1{,}024$ in the first extension and $128$ in the second}; index the copies by
$k = 1, \dots, C$. Define $O$ as the number of open cells of the maze map, $O = 46$ for the large
map.

The task is continuing, so an episode ends only at the $400$-step truncation, a copy's four
environments are reset together, and one iteration collects $T \times N = 512$ environment steps
per copy --- $128$ steps in each of the four. Define the \emph{episode phase} of iteration $i$ as
$\phi_i = ((i - 1)\,T) \bmod 400$, the point inside the episode at which that iteration's rollout
begins. It cycles with period $400 / \gcd(128, 400) = 25$ iterations, so one iteration's extrinsic
reward is the reward collected in one $128$-step slice of the episode and can read exactly zero for
every copy while copies are reaching the goal. No single iteration is therefore a score. Define a
\emph{recorded window} as a block of $200$ consecutive iterations --- eight whole turns of that
cycle, and exactly $E = 256$ complete episodes per copy --- and define $W = 98$ as the number of
windows in the run. A window covers every part of the episode the same number of times, so the
phase cancels out of its total. Define $S_{k,w}$ as the extrinsic reward copy $k$ collects over the
whole of window $w$; that total, per copy, is what the trainer records, and it is the only quantity
the metrics below are built from.

Every table cell is the mean over the configuration's completed copies plus or minus its standard
error; $N$ is how many completed copies that mean is taken over.

\begin{itemize}[leftmargin=*]
\item \textbf{Whole-run reward} (higher is better; the column the table is sorted by) --- copy
  $k$'s mean episode return over its whole run,
  $\bar{R}_k = \dfrac{1}{W E} \sum_{w=1}^{W} S_{k,w}$.
\item \claudeedit{\textbf{Max reward} (higher is better) --- the highest point the configuration's
  \emph{aggregated} curve reaches. Define $m_w$ as the mean over the configuration's copies of the
  window-$w$ episode return, $m_w = \dfrac{1}{K} \sum_{k=1}^{K} \dfrac{S_{k,w}}{E}$, where $K$ is
  the number of completed copies; the metric is $\max_{w \le W} m_w$, reported with its standard
  error over the copies at that same window. \textbf{The copies are averaged first and the maximum
  taken second.} Taking each copy's own best window and averaging those would give a larger number
  describing a run that never happened, because every copy peaks at a different moment. Read
  against whole-run and last-window reward it says how much of a configuration's best a longer
  budget gives back: a configuration whose max and last-window agree never fell, and one whose max
  is far above its last-window peaked and declined.}
\item \textbf{Last-window reward} (higher is better) --- the same mean over the final window alone,
  $\bar{R}^{\,\mathrm{last}}_k = \dfrac{S_{k,W}}{E}$, that is over copy $k$'s last $256$ episodes.
  It separates a configuration that ends its run at the goal from one that reached the goal early
  and then wandered off. It stands in for the last-$100$-episode mean a per-episode record would
  allow: this trainer records whole windows, and $256$ episodes is the shortest window whose total
  is free of the episode phase.
\item \textbf{Success rate} (higher is better) --- the fraction of the configuration's copies that
  reached the goal at least once,
  $s = \dfrac{1}{C} \sum_{k=1}^{C} \mathbf{1}\!\left[\sum_{w=1}^{W} S_{k,w} > 0\right]$. A copy that
  never enters the goal radius collects exactly zero, because reaching the goal is the only thing
  this environment pays for.
\item \textbf{Maze-cell coverage \%} (higher is better) --- define $V_k$ as the number of distinct
  open maze cells copy $k$ has entered at least once over its whole run; the metric is
  $c_k = 100 \times \dfrac{V_k}{O}$.
\item \textbf{Steps to goal} (lower is better) --- \textbf{not recorded in this run.} The metric is
  the mean, over a copy's successful episodes only, of the step inside the episode at which it first
  earns reward $1$. This trainer never reads that step back to the host: the goal step is a
  per-step quantity, and a per-step host synchronisation would slow every iteration of every one of
  the $8{,}448$ copies. Recording it needs a fixed-shape accumulator in the training state behind a
  static configuration flag, read once per window --- the pattern maze-cell coverage already uses.
  That was not added mid-campaign; the request is logged in the platform's
  \texttt{research/\allowbreak README.md}, and the column prints ``not recorded'' rather than an
  empty cell.
\end{itemize}

\begin{table}[H]
\centering
\footnotesize
\setlength{\tabcolsep}{3pt}
\renewcommand{\arraystretch}{1.15}
\resizebox{\textwidth}{!}{%
% Auto-generated by code/<YYYY-MM-DD-HH-MM>_<run-slug>-tables/make_all.py once completed records
% exist. Everything between the two markers is replaced wholesale; the caption, the label and the
% sizing stay outside them. Until then the block below holds the final column headers and one row
% saying the run has not produced a record yet.
% >>> AUTO-GENERATED TABLE START: pm-jax-run11-results
\begin{tabular}{@{}>{\raggedright\arraybackslash}p{4.4cm} r r r r r r r@{}}
\toprule
\textbf{algorithm --- best configuration} &
\textbf{\shortstack[c]{whole-run\\reward $\uparrow$}} &
\textbf{\shortstack[c]{max\\reward}} &
\textbf{\shortstack[c]{last-window\\reward}} &
\textbf{\shortstack[c]{success\\rate}} &
\textbf{\shortstack[c]{maze-cell\\coverage \%}} &
\textbf{\shortstack[c]{steps\\to goal}} &
\textbf{$N$} \\
\midrule
\multicolumn{8}{@{}l}{\textbf{\texttt{pointmaze\_\allowbreak large\_\allowbreak cont400\_\allowbreak nonoise@1}}} \\
\texttt{gt\_\allowbreak position\_\allowbreak velocity\_\allowbreak sqrt}\newline \small learning rate $10^{-3}$, $\beta = 10^{0}$ & {\boldmath $14.608 \pm 0.981$} & {\boldmath $31.417 \pm 1.916$} & {\boldmath $31.339 \pm 1.915$} & {\boldmath $0.945$} & \underline{$96.51 \pm 0.48$} & not recorded & $256$ \\
\texttt{rnd\_\allowbreak next\_\allowbreak state}\newline \small learning rate $10^{-3}$, $\beta = 10^{1}$ & \underline{$7.866 \pm 0.723$} & \underline{$17.070 \pm 1.464$} & \underline{$17.070 \pm 1.464$} & \underline{$0.926$} & {\boldmath $96.86 \pm 0.48$} & not recorded & $256$ \\
\texttt{gt\_\allowbreak position\_\allowbreak velocity\_\allowbreak linear}\newline \small learning rate $10^{-3}$, $\beta = 10^{0}$ & $2.423 \pm 0.464$ & $5.667 \pm 1.031$ & $5.654 \pm 1.026$ & $0.500$ & $80.06 \pm 1.35$ & not recorded & $256$ \\
\texttt{no\_\allowbreak exploration}\newline \small learning rate $10^{-3}$, no intrinsic weight & $0.005 \pm 0.005$ & $0.026 \pm 0.026$ & $0.019 \pm 0.019$ & $0.004$ & $24.07 \pm 0.89$ & not recorded & $256$ \\
\midrule\midrule
\texttt{rnd\_\allowbreak next\_\allowbreak state}\newline \small learning rate $10^{-3}$, $\beta = 10^{1}$ & {\boldmath $1.795 \pm 0.101$} & \underline{$25.776 \pm 0.955$} & {\boldmath $1.242 \pm 0.286$} & {\boldmath $0.999$} & {\boldmath $99.58 \pm 0.06$} & not recorded & $1{,}024$ \\
\texttt{gt\_\allowbreak position\_\allowbreak velocity\_\allowbreak sqrt}\newline \small learning rate $10^{-3}$, $\beta = 10^{0}$ & \underline{$0.658 \pm 0.024$} & {\boldmath $31.447 \pm 1.053$} & \underline{$0.185 \pm 0.118$} & {\boldmath $0.999$} & \underline{$99.07 \pm 0.06$} & not recorded & $1{,}024$ \\
\texttt{gt\_\allowbreak position\_\allowbreak velocity\_\allowbreak linear}\newline \small learning rate $10^{-3}$, $\beta = 10^{0}$ & $0.331 \pm 0.030$ & $7.181 \pm 0.587$ & $0.171 \pm 0.095$ & \underline{$0.804$} & $94.50 \pm 0.35$ & not recorded & $1{,}024$ \\
\texttt{no\_\allowbreak exploration}\newline \small learning rate $10^{-3}$, no intrinsic weight & $0.001 \pm 0.001$ & $0.088 \pm 0.088$ & $0.000 \pm 0.000$ & $0.016$ & $33.36 \pm 0.53$ & not recorded & $1{,}024$ \\
\midrule\midrule
\texttt{rnd\_\allowbreak next\_\allowbreak state}\newline \small learning rate $10^{-4}$, $\beta = 10^{2}$ & {\boldmath $52.500 \pm 3.368$} & {\boldmath $78.532 \pm 4.113$} & {\boldmath $50.754 \pm 4.851$} & {\boldmath $1.000$} & {\boldmath $99.85 \pm 0.06$} & not recorded & $128$ \\
\texttt{gt\_\allowbreak position\_\allowbreak velocity\_\allowbreak sqrt}\newline \small learning rate $10^{-4}$, $\beta = 10^{1}$ & \underline{$38.199 \pm 3.216$} & \underline{$60.835 \pm 4.816$} & \underline{$33.854 \pm 4.629$} & {\boldmath $1.000$} & {\boldmath $99.85 \pm 0.07$} & not recorded & $128$ \\
\bottomrule
\end{tabular}
% <<< AUTO-GENERATED TABLE END: pm-jax-run11-results
}
\caption{Best configuration per algorithm arm, ranked by whole-run reward (the only arrowed
column). \textbf{\sout{The double rule separates the two budgets.}\claudeedit{The two double rules
separate three runs.}} \sout{Above it are}\claudeedit{Above the first are} the four winners of the
original sweep, $256$ copies each over \sout{$10{,}035{,}200$ environment}\claudeedit{$10$M} steps
per copy; below it are
the same four configurations as the extension \claudeedit{of 2026-08-16} re-ran them, $1{,}024$
copies each over
\sout{$1{,}000{,}038{,}400$}\claudeedit{$1000$M} steps per copy\claudeedit{; below the second are
the two arms whose whole grid the extension of 2026-08-17 re-ran at $1000$M steps, each at the
configuration that grid selects, $128$ copies each. The two lower blocks share a budget and differ
only in how their configuration was chosen --- by a $10$M-step sweep, or by a $1000$M-step one}.
\claudeedit{The max-reward column is the highest point each configuration's copy-averaged curve
reaches, which is where it differs from every other column: the rest describe the end of the run or
an average over it, and that one describes its best moment.} Cells are mean $\pm$ standard error over the configuration's
completed copies; $N$ is how many completed copies, and every unit of \sout{both runs}\claudeedit{all
three runs} finished, so $N$ is
the full $256$ above the \claudeedit{first} double rule\sout{ and}\claudeedit{,} the full $1{,}024$
below it\claudeedit{, and the full $128$ below the second}. Best per metric column
\textbf{bold}, second best \underline{underlined}, computed \emph{within} each block --- the
\sout{two }blocks are different step budgets over different copy counts, so a cell of one is not
comparable with a cell of \sout{the other}\claudeedit{another}. Higher is better in every column that carries a number. Steps to goal
would be the one column where lower is better, and this trainer does not record it --- the cell says
so rather than being left empty.
\claudecomment{Budgets rounded and the run named by date, under the shared skill
\texttt{rl-experiment-writing}: the exact budgets are $10{,}035{,}200$ steps per copy for the
original sweep and $1{,}000{,}038{,}400$ for the extension, recorded in each run's
\texttt{manifest.yaml} (\texttt{compile\allowbreak\_signature.\allowbreak
env\allowbreak\_steps\allowbreak\_per\allowbreak\_copy}). The date added here is the extension
run's, and it replaces the run-folder path that the what-is-swept table used to carry.}}
\label{tab:pm-jax-run11-results}
\end{table}

% ARTIFACT 4 --- the training-curve figure. Drawn by <run>/analysis/code/make_reward_curves.py,
% which imports the same best_per_arm the results table above uses, and regenerated with it. The
% run folder's name carries underscores, so the path is wrapped in \detokenize.
% The extension replaced the curves rather than adding to them: the path now names the EXTENSION
% run folder, whose script writes the same file name into its own analysis/plots/. The original
% run's plot is untouched in the original run folder and is its record of the shorter budget.
\begin{figure}[H]
\centering
\includegraphics[width=\textwidth]{\detokenize{../runs/2026-08-17-10-45_jax-ppo_pointmaze-large-nonoise_full-batch_rollout-128_envs-per-copy-4_steps-per-copy-1000M_bonus-rnd-next-state-visit-count-sqrt_learning-rate-1e-3-1e-4-1e-5_intrinsic-weight-1e-5-to-1e5_copies-per-cell-128_extension-2-of-first-batch/analysis/plots/pm_jax_run11_reward_curves.pdf}}
\caption{Mean episode return against environment steps for each algorithm arm's best
configuration, chosen by the same score that ranks \Cref{tab:pm-jax-run11-results}. \textbf{The
budget shown is the extension's}: every curve runs to
\sout{$1{,}000{,}038{,}400$ environment}\claudeedit{$1000$M} steps per
copy\sout{ and is the mean over that configuration's $1{,}024$ copies}, with the band plus or minus
one standard error over \sout{them}\claudeedit{its copies}; each point is one recorded window of
$256$ episodes per copy. \claudeedit{\textbf{Line style says which sweep chose the configuration,
not which budget it ran at.} The four \textbf{dashed} curves are the configurations the $10$M-step
sweep picked, as the first extension ran them over $1{,}024$ copies; the two \textbf{solid} ones
are the configurations the second extension's own $1000$M-step sweep picks for the two arms it
re-ran, over $128$ copies. Both sets run to the same $1000$M steps, so the gap between them is
what choosing at the wrong budget costs.}
\textbf{The step axis is logarithmic}, which the original sweep's figure did not need: \sout{every
arm reaches its peak inside the first $15$ million steps and then decays over the remaining $985$
million}\claudeedit{the dashed curves peak inside the first $15$ million steps and the solid ones
near $200$ million, two orders of magnitude apart}, so on a linear axis the whole \sout{rise would
occupy}\claudeedit{dashed set would occupy} the first per cent of the width. The
original sweep's $10$-million-step curves over $256$ copies are not drawn beside them --- two
budgets of one arm in one panel read as two algorithms --- and stay in the original run folder's
own plot. Solid black is the best random-network-distillation configuration\claudeedit{ of the
$1000$M-step sweep}, the reference the
other arms are measured against; every other arm keeps the colour \sout{and line style }it had in the
original run. The legend names each winner's learning rate $\alpha$ and, where the arm has a
bonus, its intrinsic weight $\beta$. Drawn by
\texttt{analysis/\allowbreak code/\allowbreak make\allowbreak\_reward\allowbreak\_curves.py} in the
\sout{extension}\claudeedit{second extension's} run folder.
\claudecomment{The budget is written round under the shared skill
\texttt{rl-experiment-writing}; the exact $1{,}000{,}038{,}400$ steps per copy is in the extension
run's \texttt{manifest.yaml}. The figure file's own path is unchanged --- it is the
\texttt{\textbackslash includegraphics} argument the build needs, not prose.}}
\label{fig:pm-jax-run11-curves}
\end{figure}

\paragraph{Findings.}
\begin{enumerate}[label=\textbf{F\arabic*},leftmargin=*,labelsep=0.6em]
\item \textbf{H1 holds wherever reward can separate the arms.} At the learning rate $10^{-3}$
  \texttt{rnd\allowbreak\_next\allowbreak\_state} at its best weight reaches a whole-run reward of
  $7.866 \pm 0.723$ against $0.005 \pm 0.005$ for \texttt{no\allowbreak\_exploration}, and $0.926$
  of its copies reach the goal at least once against $0.004$. At $10^{-4}$ the gap is
  $0.174 \pm 0.031$ against $0.000 \pm 0.000$. At $10^{-5}$ neither arm ever reaches the goal, so
  reward is exactly zero on both sides and cannot separate them --- but coverage still can:
  $42.58 \pm 0.82$ per cent of the maze against $29.69 \pm 0.60$. The bonus is working there; the
  sparse reward is simply not yet reporting it.
\item \textbf{H2 is refuted as stated: only one of the two oracle bonuses bounds distillation from
  above.} The $1/\sqrt{n}$ arm does, and by a wide margin --- $14.608 \pm 0.981$ against
  $7.866 \pm 0.723$. The $1/n$ arm does not: $2.423 \pm 0.464$, well below distillation. And even the
  winning oracle arm does not bound distillation on coverage: $96.51 \pm 0.48$ against
  $96.86 \pm 0.48$, a difference smaller than either standard error. Having the true novelty signal
  is therefore not sufficient on its own; how fast the bonus decays with the count decides whether
  the signal is worth more than a predictor's error.
\item \textbf{H3 is refuted: the two decays do not trade coverage against reward, one is simply
  better.} The prediction was that $1/\sqrt{n}$ would take the coverage and $1/n$ the reward. The
  $1/\sqrt{n}$ arm takes both, $96.51 \pm 0.48$ against $80.06 \pm 1.35$ on coverage and
  $14.608 \pm 0.981$ against $2.423 \pm 0.464$ on reward. On this maze $1/n$ stops paying for a
  state too early to finish exploring it.
\item \textbf{H4 half-holds: there is one interior peak and a sharp decline above it, but no rise
  below it.} At $10^{-3}$ the distillation arm goes $2.15$ at $\beta = 10^{0}$, $7.87$ at $10^{1}$,
  $2.34$ at $10^{2}$, $0.22$ at $10^{3}$; the $1/\sqrt{n}$ arm goes $1.68$, $14.61$, $7.30$, $0.39$
  over $\beta = 10^{-1}$ to $10^{2}$. Below the peak, though, every weight at or under $10^{-2}$
  sits on the zero floor in every arm, so the ordering there is noise rather than a rise. The
  decline above the peak behaves exactly as the hypothesis said it would: coverage stays saturated
  at $97$ to $99\%$ while reward collapses, which is the agent still chasing novelty after the goal
  is known.
\item \textbf{A finding outside the hypotheses: the learning rate decides more than the bonus does.}
  Only $10^{-3}$ produces a configuration that solves the maze. At $10^{-4}$ the best arm reaches
  $0.174$, and at $10^{-5}$ every one of the four arms reaches exactly zero. The three rates are
  therefore not three comparable settings of one experiment --- one of them works, and the next run
  should sweep upward from $10^{-3}$ rather than downward from it.
\item \claudeedit{\textbf{F5 is refuted by the second extension: the learning rate that wins at
  $10$M is not the one that wins at $1000$M, and neither arm's winning configuration survives the
  change of budget.} Given its whole grid at the long budget, random network distillation wins at
  learning rate $10^{-4}$ and $\beta = 10^{2}$ with $52.500 \pm 3.368$, and the oracle $1/\sqrt{n}$
  bonus at $10^{-4}$ and $\beta = 10^{1}$ with $38.199 \pm 3.216$. The configurations the $10$M
  sweep picked --- both at $10^{-3}$, at $\beta = 10^{1}$ and $\beta = 10^{0}$ --- score
  $1.795 \pm 0.101$ and $0.658 \pm 0.024$ over the same $1000$M steps: a factor of $29$ and $58$
  worse than the arm's own best. Both arms move the same way, one step down in learning rate and
  one step up in intrinsic weight, so this is a property of the budget rather than of either
  bonus. A short sweep chooses what climbs fastest, and what climbs fastest is not what climbs
  furthest.}
\item \claudeedit{\textbf{The ordering between the two arms is unchanged, and it is now measured at
  each arm's own best.} Distillation leads the oracle $1/\sqrt{n}$ bonus $52.500 \pm 3.368$ against
  $38.199 \pm 3.216$, a gap of about four standard errors. That is the same direction the first
  extension found and the opposite of the $10$M sweep's, and it no longer depends on the two arms
  having been given configurations chosen for a hundred-times-shorter run. Neither arm is separated
  on the other two metrics: both reach the goal in every one of their $128$ copies and both cover
  $99.85$ per cent of the maze, so at this budget reward is the only metric with anything left to
  say.}
\end{enumerate}

\paragraph{\claudeedit{Curve shape: does every configuration rise and then come down?}}
\claudeedit{\textbf{No.} Aggregated over a configuration's $128$ copies and classified by
\texttt{analysis/\allowbreak code/\allowbreak curve\allowbreak\_shape.py}, $20$ of the $66$
configurations ($30$ per cent) rise and then fall; $16$ ($24$ per cent) never reach the goal at all,
$14$ ($21$ per cent) are still rising at $1000$M steps, $13$ ($20$ per cent) decline from the start,
and $3$ ($5$ per cent) rise and then flatten. The classifier's label is unchanged under all nine of
its threshold variants for $53$ of the $66$, and the instability is confined to configurations that
score badly: every one of the top ten by whole-run reward, from $52.500$ down to $11.631$, carries
the same label under all nine variants, while the thirteen unstable ones all score $10.391$ or
below.}

\claudeedit{Because neither arm's best configuration is one of the exceptions ---
\Cref{tab:pm-jax-run11-results} and \Cref{fig:pm-jax-run11-curves} carry two curves that both rise
and fall --- the best exception of each arm is a configuration those artifacts do not show, and
\Cref{tab:pm-jax-run11-shapes} is where they appear.}

\claudeedit{What decides the shape is the learning rate, and \Cref{fig:pm-jax-run11-shapes} is
where that is visible. It draws each arm at each of the three swept learning rates, at the
intrinsic weight whose curve reaches the highest point, and the three rates behave in three
different ways at the same budget: $\alpha = 10^{-3}$ peaks inside the first $15$M steps and then
collapses --- the oracle arm to a mean episode return of zero --- $\alpha = 10^{-4}$ peaks near
$200$M an order of magnitude higher and gives about half of it back, and $\alpha = 10^{-5}$ is
still climbing in the last window and has the lowest peak of the three. The two arms track each
other closely within a rate, which is the evidence that the rate and not the bonus is what governs
when a run peaks.}

\begin{table}[H]
\centering
\footnotesize
\setlength{\tabcolsep}{3.5pt}
\renewcommand{\arraystretch}{1.15}
\begin{tabular}{@{}>{\raggedright\arraybackslash}p{3.4cm} r r >{\raggedright\arraybackslash}p{1.9cm} r r r@{}}
\toprule
\textbf{algorithm} & \textbf{\shortstack[c]{learning\\rate}} &
\textbf{\shortstack[c]{intrinsic\\weight}} & \textbf{aggregate shape} &
\textbf{\shortstack[c]{whole-run\\reward $\uparrow$}} & \textbf{\shortstack[c]{rank of\\$66$}} &
\textbf{\shortstack[c]{copies of $128$\\sharing it}} \\
\midrule
\texttt{rnd\_\allowbreak next\_\allowbreak state} & $10^{-5}$ & $10^{1}$ & still rising &
  $27.778$ & $6$ & $55$ \\
\texttt{gt\_\allowbreak position\_\allowbreak velocity\_\allowbreak sqrt} & $10^{-4}$ & $10^{2}$ &
  rise then plateau & $10.391$ & $11$ & $6$ \\
\midrule
\multicolumn{7}{@{}l}{\textit{for comparison, the same arms' best configurations, which do rise then
  fall}} \\
\texttt{rnd\_\allowbreak next\_\allowbreak state} & $10^{-4}$ & $10^{2}$ & rise then fall &
  $52.500$ & $1$ & $76$ \\
\texttt{gt\_\allowbreak position\_\allowbreak velocity\_\allowbreak sqrt} & $10^{-4}$ & $10^{1}$ &
  rise then fall & $38.199$ & $3$ & $84$ \\
\bottomrule
\end{tabular}
\caption{\claudeedit{The best configuration of each arm that does \emph{not} rise then fall, with
each arm's overall best beneath the rule for comparison. Rank is over all $66$ configurations of
this sweep by whole-run reward. \textbf{The last column is the caveat this table exists to carry:}
the shape of the mean over $128$ copies need not be the shape of any copy, and it often is not ---
$76$ and $84$ of the $128$ copies share the aggregate label of the two best configurations, but
only $6$ of $128$ share the oracle arm's rise-then-plateau, where $90$ of its copies rise and fall
on their own. The oracle row's label is also the one label here that moves under the threshold
variants, reading \emph{still rising} under two of the nine; no variant calls it rise-then-fall, so
the answer to the question is unaffected, but which non-falling shape it is is not settled.}}
\label{tab:pm-jax-run11-shapes}
\end{table}

\begin{figure}[H]
\centering
\includegraphics[width=\textwidth]{\detokenize{../runs/2026-08-17-10-45_jax-ppo_pointmaze-large-nonoise_full-batch_rollout-128_envs-per-copy-4_steps-per-copy-1000M_bonus-rnd-next-state-visit-count-sqrt_learning-rate-1e-3-1e-4-1e-5_intrinsic-weight-1e-5-to-1e5_copies-per-cell-128_extension-2-of-first-batch/analysis/plots/pm_jax_run11_curve_shapes.pdf}}
\caption{\claudeedit{Each arm at each of the three swept learning rates, six curves in all. The
curve drawn for a pair is the intrinsic weight whose aggregated curve reaches the highest point,
and that point is marked on the curve itself and labelled beside the marker, so the best a
configuration ever reaches can be read off without matching a line to a legend entry. \textbf{The
maximum is taken after averaging, not before}: the curve is the mean over the configuration's $128$
copies window by window, and the mark is the largest value that mean reaches. The mean of each
copy's own best window would be larger and would describe no run that ever happened, since every
copy peaks at a different moment. Colour is the arm; line style and transparency are the learning
rate, dotted and faint for $10^{-3}$, solid for $10^{-4}$, dashed for $10^{-5}$. The band is plus
or minus one standard error over the $128$ copies. Note that the weight selected here need not be
the one \Cref{tab:pm-jax-run11-results} names, because that table ranks on whole-run reward and
this figure ranks on the peak: for distillation at $\alpha = 10^{-4}$ the two criteria choose
$\beta = 10^{3}$ and $\beta = 10^{2}$ respectively. Drawn by
\texttt{analysis/\allowbreak code/\allowbreak make\allowbreak\_shape\allowbreak\_curves.py} in the
second extension's run folder.}}
\label{fig:pm-jax-run11-shapes}
\end{figure}

\noindent\claudecomment{This table and this figure are an addition the extension convention of
\texttt{rnd-experiment-tex-track} would not normally allow, and they are here because the questions
they answer were asked of this run directly and neither answer is visible in
\Cref{tab:pm-jax-run11-results} or \Cref{fig:pm-jax-run11-curves}. The table's question was whether
every configuration rises then falls: it does not, and because both of those artifacts carry each
arm's best configuration and both of those rise and fall, the exceptions appear nowhere else. Had
either arm's best been an exception, the specification asked for no separate table and none would
have been added. The figure answers a second question asked at the same time --- each arm's best
configuration at each of the three learning rates, six curves, the best defined as the highest
point the copy-averaged curve reaches --- and it replaced an earlier version of this figure that
drew the table's four configurations instead. Sources: the shapes and their sensitivity from
\texttt{analysis/\allowbreak shape\allowbreak\_classification.json} in the second extension's run
folder, the per-copy agreement from
\texttt{analysis/\allowbreak code/\allowbreak copy\allowbreak\_agreement.py}, and every curve
through the same \texttt{code/\allowbreak aggregate.\allowbreak curve\allowbreak\_of} the results
table uses, gated identical against a one-pass reader before the classification ran.}

\claudecomment{A \textbf{Training cost} paragraph and a table of what the run cost on each
card --- seconds per iteration, total and per-copy step rates, hours per million steps, build
and compile time, wall clock --- stood here and were deleted, not struck through: rule~3 of the
shared skill \texttt{rl-experiment-writing} keeps infrastructure telemetry out of a development
document unless the reader asks for it. The measurements are unchanged and still live in the run
folder's \texttt{summary.json} and \texttt{infra\allowbreak\_history.md}, and the extension run's
\texttt{code/\allowbreak submission\allowbreak\_plan.json} holds the per-card plan. The table
carried no number that any other part of this document cites, and its label
\texttt{tab:pm-jax-run11-cost} was referenced only from inside the deleted paragraph.}

\end{document}
